\section{Ablation Results}
\label{sec:results}

\textbf{Enabling all three modules raised end-to-end success from 52.1\%
to 100.0\% --- an absolute gain of $+47.9$ percentage points (95\%
bootstrap CI $[+39.6, +56.3]$), a relative reduction of 100\% of the
baseline failure rate, on the 144-cell combined-stress grid at $N=21$,
$a=19$, $r=6$.} All 69 baseline failures were converted, and no trial
regressed ($b=0$, $c=69$; exact McNemar $p=3.4\times10^{-21}$).

This is a within-regime result and should be read as one: it is the
comparison the ablation was designed to make, and Secs.~\ref{sec:heldout}
and \ref{sec:matched} each qualify it materially.

Per-arm figures for this grid appear as the first block of
Table~\ref{tab:heldout}, alongside the held-out instances, so that the
canonical and replication results can be read against each other rather
than in separate tables. Module B (candidate coverage) and module D (shot
stopping) each recover roughly half the baseline failure rate on their
own, at $+25.0$ and $+24.3$ percentage points. Module C (beam width)
moves the success rate by only $+4.9$ points on seven discordant trials
and does not survive Holm correction at the confirmatory family size
(\SItables{}); it should be read as a null result rather than a
small positive one. The planned B-vs-C contrast confirms the asymmetry:
B outperforms C by $20.1$ points (95\% CI $[+12.5,+29.2]$, $(b,c)=(8,37)$
counting C as the reference, $p=1.5\times10^{-5}$), and the same contrast
reproduces on the held-out $(N{=}21,a{=}2,r{=}6)$ instance at $+21.5$
points ($p=3.1\times10^{-6}$).



\paragraph{The modules are not interchangeable.} B at 77.1\% and C at
56.9\% do not compose to Full's 100\%, consistent with the attribution of
Sec.~\ref{sec:attribution}: 20 of the 69 hardest failures needed candidate
budget \emph{and} beam width relaxed together, and neither axis alone
reaches them. D's effect size closely matching B's likewise corroborates
the shot-limitation attribution (48 of 69) --- a prediction the design
made in advance rather than a post-hoc reading.

\paragraph{Individual modules can regress individual trials.} Module B
turned two baseline successes into failures, and module D six. These are
not measurement noise; they are the expected consequence of a point made
in the companion article~\cite{methodpaper}. Relaxing a budget changes which
candidates are retained and how the resulting paths rank, and the
continued-fraction stage is not monotone in the candidate set --- a
different phase can be tried first and fail where the baseline's would
have succeeded. Both regressions are absorbed in the Full arm, which
regresses no trial on any of the four instances, but a practitioner
enabling B or D \emph{alone} should expect a small rate of such
reversals; on $(N{=}15,a{=}2,r{=}4)$ module D alone breaks exactly as
many trials as it fixes (Sec.~\ref{sec:heldout}). We flag this because
the structural monotonicity result invites the opposite conclusion.
